\begin{statement}{C.2.UN.57}{Proved}{The exact marked-cycle criterion without an algebraic-dimension restriction}
Specify a connected oriented cycle of marked surfaces, each either a
diagonal primary or secondary Hopf surface with its two distinct axes,
or $\PP(\OO_E\oplus L)$ with $L\in\operatorname{Pic}^0(E)$ and its
zero and infinity sections. Specify the elliptic gluing maps and
duality of the actual axis or section normal lines. Components are
smooth embedded components, except that a single ruled normalization
with its two sections identified is also allowed.

The actual rank-three complement-lattice product $M$ is the identity
precisely for those specified supports which admit a smooth-total
two-torus smoothing. Every such support occurs as a full reduced fibre
of a smooth compact complex threefold over $\PP^1$. The construction
retains each component and every specified elliptic gluing map. Its
only other singular fibre is the complex conjugate marked support,
and every reduced fibre support has Euler number zero. No restriction
on the algebraic dimension of the diagonal Hopf components is imposed.
The construction chooses normal-product constants.
\end{statement}

\begin{proof}
\par\medskip\noindent\textbf{The actual complement maps and the local criterion.}
For a diagonal Hopf component the exponential preimage of its full
diagonal deck group is an actual rank-three lattice $\Lambda_i$ in
$\C^2$. Its complement of both axes is $\C^2/\Lambda_i$. This
description retains the finite secondary group, including its two
angular characters. Both small normal meridians are primitive: a
smaller angular lattice vector would give a nonidentity finite deck
element fixing a nonzero point of the other axis, contrary to the
freeness of this Hopf quotient. For a ruled component the complement
is exactly $L^\times$; its degree-zero line-bundle factors give the
rank-three lattice and primitive normal meridian in (C.MX.3).

An elliptic gluing map lifts to an affine map of universal covers.
A dual-normal line isomorphism lifts to a nowhere-zero function;
the derivative of its logarithm is periodic under the two elliptic
periods, hence constant. Its logarithm is therefore affine.
Consequently the actual normal-product neck has exactly the affine
logarithmic maps of (C.HC.5) and (C.MX.2), including their constants.
These statements use the displayed axes and sections; they use no
absence of further curves on a component.

The component cores and the $xy=t$ necks give the actual fundamental
group
\[
                         \pi_1(F_t)=\Lambda\rtimes_M\Z.
                                                        \tag{C.AD.1}
\]
For an arbitrary specified smoothing, this is the same comparison
as in UN46 and UN51: its canonical divisor sections and a smooth
tubular retraction identify the component cores with their actual
normal circle bundles, and its two divisor sections identify each
double collar with $xy=t$. The changes of these identifications
give the specified elliptic and normal-line maps on fundamental
groups. Thus (C.AD.1) also holds for the actual smoothing, without
asserting that its holomorphic neighbourhood is linear. The
abelian group of a two-torus forces $M=I$.

Conversely the actual flat normal lines and affine necks give the
proper normal-product plumbing of those proofs. Their argument
requires only the entire line-bundle axis charts, disjoint marked
curves and the displayed constant factors, all present here.
When $M=I$, the common logarithmic coframe glues on every nearby
compact fibre. Its two dual holomorphic fields commute and are
complete. Their $\C^2$ action has discrete cocompact stabilizer
and identifies that fibre with a complex two-torus. This proves
the full local equivalence in the present generality.

\par\medskip\noindent\textbf{The positive total normal period.}
Assume $M=I$ and identify the complement groups with
$G=\C^2/\Lambda$ by the actual linear edge maps. Write $n_i$
for the outgoing normal meridian at vertex $i$; the incoming
one is $-n_{i-1}$. At a Hopf vertex these two vectors are
complex-linearly independent, and its actual logarithmic frame is
\[
       P_i=\left(\frac{n_i}{2\pi i},
                         -\frac{n_{i-1}}{2\pi i}\right).
                                                        \tag{C.AD.2}
\]
At a ruled vertex $n_i=n_{i-1}$, and its separate frame is the
one in (C.MX.10).

Put $W=\operatorname{span}_{\R}\Lambda$ and choose a nonzero
real linear functional $\ell$ with kernel $W$. A contracting
Hopf period with logarithms $\log\alpha_i,\log\beta_i$ lies
in $W$. Applying $\ell$ to (C.AD.2) times that period gives
\[
 \log|\alpha_i|\,\ell(i n_i)
       =\log|\beta_i|\,\ell(i n_{i-1}).                 \tag{C.AD.3}
\]
Both logarithms are negative. Neither transverse value can
vanish: their simultaneous vanishing would put both complex
normal lines, hence all $\C^2$, inside the real three-plane
$W$. Thus they have the same nonzero sign. At a ruled vertex
the two values coincide. In an all-ruled cycle their common
value is nonzero because the actual elliptic quotient satisfies
$\C n\cap W=\R n$. Hence all values around every present
cycle have the same nonzero sign.

The exact fourth period is
\[
 b+\frac{\log t}{2\pi i}u,\qquad
             u=-\sum_i n_i\in\Lambda,
                                                        \tag{C.AD.4}
\]
where $b$ is the sum of the actual affine edge constants.
Equation (C.AD.3) gives $\ell(iu)\ne0$. Write $u=d p$ with
$p$ primitive and $d>0$. The actual complex linear coordinate
change of (C.TL.1) therefore gives
\[
 \begin{gathered}
 W=\C\times\R,\qquad p=(0,1),\\
 \Lambda=\langle (a_1,c_1),(a_2,c_2),p\rangle_{\Z},
 \quad c_1,c_2\in\R,\quad
 C=\operatorname{Im}(\overline{a_1}a_2)>0.
 \end{gathered}                                        \tag{C.AD.5}
\]
Every $n_i$ has nonzero second coordinate, by the transverse
calculation. Changing one normal-product constant changes
its affine edge constant by an arbitrary multiple of $n_i$.
Choose it to make $b=(\sigma,i)$. This changes no elliptic
gluing map. In the all-ruled case it leaves the full horizontal
translation $\sigma$ unchanged, exactly as in UN55.

\par\medskip\noindent\textbf{The entire punctured sphere of actual tori.}
Use $V_z$ from (C.TL.3), and set
\[
 \begin{gathered}
 \Lambda(z)=V_z\Lambda,\qquad
 B(z)=V_z(b)+\frac{d}{2\pi i}\log z\,V_z(p),\\
                   T_z=\C^2/(\Lambda(z)+\Z B(z)).
 \end{gathered}                                        \tag{C.AD.6}
\]
These are holomorphic period sections for every $z\in\C^*$.
The change of a logarithm adds the integral period $dV_z(p)$.
The exact real determinant proved in (C.TL.5)--(C.TL.6), with
$N$ replaced by the positive integer $d$, is
\[
 C(1+|z|^2)\left(1+|z|^2+
          \frac{d}{2\pi}(|z|^2-1)\log|z|\right)>0.
                                                        \tag{C.AD.7}
\]
Thus (C.AD.6) gives an actual smooth proper family of two-tori
on all $\C^*$, by the locally uniformly bounded lattice
matrices and parallelepipeds in UN55. No elliptic quotient
of this entire middle family is required.

\par\medskip\noindent\textbf{The exact moving components at zero.}
Transport every lattice vector by $V_z$, retaining its integral
coefficients. At a Hopf vertex put
\[
 P_i(z)=\left(\frac{V_z(n_i)}{2\pi i},
                       -\frac{V_z(n_{i-1})}{2\pi i}\right),
 \qquad \Lambda_i(z)=P_i(z)^{-1}\Lambda(z).
                                                        \tag{C.AD.8}
\]
The frame is invertible near zero. Its two distinguished
angular periods are literally $(2\pi i,0)$ and $(0,2\pi i)$;
the finite angular enlargement is retained. The finite quotient characters
are unchanged, since all their integer relations were
transported. An original contracting generator becomes a
holomorphic pair of logarithms, still with negative real
parts on a smaller disc. Exponentiation therefore gives an
actual smooth family of the marked primary or secondary
diagonal Hopf quotients. Its central member is the original
marked surface, including every finite character.

At a ruled vertex use its primitive normal vector and a
transverse integral lattice lift, as in (C.MX.10). Transporting
both gives the actual moving elliptic base, degree-zero line
and projective completion. This is valid on a smaller disc
because its two base periods remain real-linearly independent.
Its two marked sections are retained.

If $b_i$ is the original common affine edge constant after
the one permitted normal-product change, choose its moving
value to be $V_z(b_i)$. The common lattice and the paired
normal meridians give the required linear edge maps. Their
affine constants induce actual elliptic translations and
normal-line isomorphisms, extending over the marked curves.
The sum of the edge translations is $V_z(b)$; the sum of
their logarithmic coefficients is $dV_z(p)/(2\pi i)$.
Thus the proper moving plumbing with equation $xy=z$ has
exactly the periods (C.AD.6). Its logarithmic chart maps
identify its entire punctured family with that family:
each collar change is the stated affine map, every branch
change is an integral period, and the complete circuit is
exactly $B(z)$. The inverse is the closed-coframe inverse
already used in the local criterion. These maps and their
inverses vary holomorphically with $z$.

\par\medskip\noindent\textbf{The conjugate Hopf and ruled filling at infinity.}
Put $s=z^{-1}$ and divide the common vector frame by $z$.
Write
\[
 W_s(w)=J\overline w+s w,\qquad
              J=\begin{pmatrix}0&1\\-1&0\end{pmatrix}.
                                                        \tag{C.AD.9}
\]
The rank-three periods become $W_s\Lambda$, and the fourth
period is
\[
                 W_s(b)-\frac{d}{2\pi i}\log s\,W_s(p).
                                                        \tag{C.AD.10}
\]
Choose the outgoing normal at vertex $i$ to be
$n_i^\infty(s)=-W_s(n_i)$. Its incoming normal is
$-n_{i-1}^\infty(s)$ and its total normal period is
$-dW_s(p)$, exactly as required by (C.AD.10).

For a Hopf vertex the corresponding frame satisfies the
following exact identities at zero:
\[
 \begin{gathered}
 P_i^\infty(0)=J\overline{P_i},\\
 (P_i^\infty(0))^{-1}W_0\Lambda
                         =\overline{\Lambda_i}.
 \end{gathered}                                        \tag{C.AD.11}
\]
The signs follow from $\overline{2\pi i}=-2\pi i$.
Consequently the actual contracting logarithms are the
complex conjugates of the original ones, whose real parts
are still negative. The finite angular characters are
their conjugates and remain free. For small $s$, the same
holomorphic lattice construction as (C.AD.8) gives the
smooth marked Hopf family through this conjugate surface.

For a ruled vertex the separate frame transforms by the
same $J$ and conjugation, with its angular normal changed
by the displayed minus sign. The resulting central ruled
surface is exactly the conjugate marked projective bundle;
the actual line-bundle lattice proves this identification.
Choose the affine constant on edge $i$ to be $W_s(b_i)$.
At zero its component-coordinate expression is the complex
conjugate of the original expression, so it induces the
conjugate elliptic gluing map. Its logarithmic coefficient
is $W_s(n_i)/(2\pi i)$, and their sum is exactly the one
in (C.AD.10).

The same proper normal-product plumbing, now with $xy=s$,
therefore has central support equal to the specified
complex conjugate support. Its actual punctured-family
logarithmic map identifies it with (C.AD.10). This proves
both the component and the entire-family comparisons at
infinity; complex conjugation is used only in the fixed
central coefficients of holomorphic period sections.

Finally glue these two proper fillings to (C.AD.6) on
their punctured discs. The exact frame change (C.AD.9)
and the two logarithmic biholomorphisms give the transition
maps. The finite closed-base refinement argument of UN55
proves the resulting total space Hausdorff and proper
over $\PP^1$, hence compact. Its displayed local models
are smooth. All middle fibres are smooth two-tori.
Every marked component and double elliptic curve has
Euler number zero, so both reduced singular supports
have Euler number zero as well.

The symbolic real determinant and the two conjugated
Hopf-frame identities are checked in
\path{code/root_all_diagonal_cycles.py} and its output.
\end{proof}


% C.2.UN.53 from part_C2_all_hopf_stars.tex

